\chapter{The Julia you need, starting from the beginning}
\label{chap:julia}

\section{What is a program?}
A program is a written set of instructions for a computer. Julia is the language in which we write those instructions. A Julia file is an ordinary text file whose name usually ends in \code{.jl}. You can open it, read it, change it, and run it. Reading a file does not run it.

Our program will not contain a tiny picture of every organism. It will store a number: how many individuals currently carry the focal allele. That is enough because our rules treat individuals of the same type identically. This is a deliberate simplification of the biology, not a limitation of the computer.

\section{Two places where you can type commands}
The \emph{terminal} is an application that accepts commands for your computer. On a Mac it may show a prompt ending with a dollar sign or a percent sign. Do not copy the prompt itself. The Julia \emph{REPL} is an interactive session that accepts Julia expressions. Its prompt is \code{julia>}. REPL means read, evaluate, print, loop: it reads your instruction, calculates the result, prints it, and waits for another instruction.

\begin{workedexample}
In the terminal, type \code{julia}. When you see \code{julia>}, type \code{2 + 3}. Julia prints \code{5}. Type \code{exit()} to leave Julia and return to the terminal. A terminal command such as \code{julia scripts/demo.jl} belongs in the terminal, not after the \code{julia>} prompt.
\end{workedexample}

The commands in this guide assume that the terminal's current directory is the project folder. A directory is a folder. The \code{cd} command changes directory. If the folder path contains spaces, put quotation marks around the whole path. To avoid needing to guess your download location, you can type \code{cd } with a trailing space, drag the extracted project folder into the terminal, and press Return.

\begin{lstlisting}[language=bash]
julia --version
julia --project=. scripts/demo.jl
\end{lstlisting}

The first command prints the installed Julia version. The second starts Julia, uses the project in the current directory (the dot means ``here''), and runs a file. The text \code{--project=.} is an option to Julia, not a line of Julia code. You do not need to type \code{julia>} anywhere in either command.

\section{Names store values}
\begin{lstlisting}
N = 5
i = 2
s = 0.5
x = i / N
\end{lstlisting}

Read the first line as ``store the number five under the name N.'' The equals sign assigns a value. It does not ask a question. The second line stores the allele count, the third stores the selective advantage, and the fourth calculates the frequency. The slash means division. Julia prints \code{0.4} if you ask for \code{x}.

An integer is a whole number such as \code{5}. A floating-point number represents a number with a fractional part, such as \code{0.5}. Most such numbers are stored with finite precision. This becomes relevant much later when we calculate probabilities very close to zero. For now, ordinary arithmetic works as you expect.

\begin{center}
\begin{tabular}{ll}
\toprule
Julia expression & Meaning\\
\midrule
\code{1 + s} & add one and the advantage\\
\code{r * i} & multiply fitness and count\\
\code{N - i} & subtract the focal count from the total\\
\code{i / N} & divide count by total\\
\code{r^2} & square the fitness\\
\code{(1 + s) * i} & do the addition before multiplying\\
\bottomrule
\end{tabular}
\end{center}

Multiplication needs a star. Write \code{r * i}, not \code{ri}. In this project a name such as \code{total_weight} is one name, not a subtraction. The underscore helps people read it.

\section{Comments explain; they do not calculate}
A hash sign starts a comment. Julia ignores everything after it on that line.
\begin{lstlisting}
x = i / N  # Frequency is count divided by population size.
\end{lstlisting}
Comments tell a reader what a calculation means. They do not make a wrong calculation correct. We will check the calculations separately.

\section{Functions are reusable instructions}
\begin{lstlisting}
function allele_frequency(i, N)
    x = i / N
    return x
end

allele_frequency(2, 5)
\end{lstlisting}

\begin{enumerate}
\item \code{function} begins a reusable instruction. \code{allele_frequency} is the name we choose for it.
\item Inside parentheses, \code{i} and \code{N} name the inputs. Inputs are also called arguments.
\item The next line performs the calculation using those inputs.
\item \code{return x} sends the answer back to the place that called the function.
\item \code{end} closes the definition. Defining a function does not yet carry out the calculation for a particular population.
\item The final line \emph{calls} the function: it supplies two and five as the inputs. Its answer is \code{0.4}.
\end{enumerate}

The order of arguments matters. \code{allele_frequency(5, 2)} asks for a different calculation. In our supplied functions the population size comes first; always check the function's documentation rather than guessing from this introductory example.

Names created inside a function normally belong to that function. A variable called \code{i} in one function does not automatically alter an unrelated \code{i} elsewhere. This helps keep a program understandable.

\section{Ask questions with comparisons}
\begin{lstlisting}
i == 0
i == N
i < N
0 <= i <= N
\end{lstlisting}

Each expression returns \code{true} or \code{false}. Two equals signs ask whether values are equal. One equals sign assigns a value. Confusing these is a common early mistake. The last expression checks whether the count is between zero and the population size, including both endpoints.

\begin{lstlisting}
if i == 0
    println("The allele has been lost.")
elseif i == N
    println("The allele has fixed.")
else
    println("Both types are still present.")
end
\end{lstlisting}

Julia tests the first question. If it is true, it runs the first message. Otherwise it tries the second question. If neither is true, it runs the \code{else} branch. Only one message is printed. \code{println} displays text and then starts a new output line. Quotation marks mark text; the messages do not affect the allele count.

\section{Loops repeat instructions}
\begin{lstlisting}
for repetition in 1:3
    println(repetition)
end
\end{lstlisting}

The colon creates the sequence one, two, three. The loop takes each number in turn, stores it in \code{repetition}, and runs its body. The output is three lines containing \code{1}, \code{2}, and \code{3}. The name \code{repetition} is descriptive; it is not a special Julia word.

A \code{while} loop repeats while a condition is true. In the model, \code{while 0 < i < N} means ``keep reproducing while both types are present.'' If an event does not change \code{i}, the loop still carried out an event. It must test the condition and try another event. A simulation can therefore take many iterations.

\section{Arrays store several values in order}
\begin{lstlisting}
counts = [2, 3, 3, 2]
counts[1]
push!(counts, 1)
length(counts)
\end{lstlisting}

Square brackets create an array, an ordered collection of values. In Julia the first entry is numbered one. Therefore \code{counts[1]} is two. The command \code{push!} adds a new entry to the end. Our array becomes \code{[2, 3, 3, 2, 1]}, and its length is five. A trailing exclamation mark is a Julia naming convention warning that a function changes something supplied to it. It is part of the function's name.

If the initial state is recorded, then after four reproduction events there are five recorded counts. Entry one is the state at event zero. Keep this one-entry difference in mind when reading a trajectory file.

\section{Random numbers select events}
\begin{lstlisting}
using Random
rng = MersenneTwister(20261027)
u = rand(rng)
\end{lstlisting}

\code{Random} is a standard Julia library. \code{using Random} makes its tools available. \code{MersenneTwister} creates a pseudorandom-number generator: a mathematical procedure producing numbers that behave like random draws for our simulation. The integer in parentheses is its seed, a starting setting. \code{rand(rng)} takes the next number from that generator. Here the result is at least zero and less than one.

Imagine a line from zero to one divided into three sections of lengths \(0.3\), \(0.2\), and \(0.5\). A uniform draw lands in each section with the corresponding probability. We can therefore use its position to choose an increase, decrease, or no change.

Setting a seed once makes a run reproducible in the recorded software environment. Setting the same seed before every repetition would repeatedly restart the same random sequence, which does \emph{not} produce independent repetitions. We keep drawing from the generator instead.

The demonstration uses some chosen numbers so that you can check the arithmetic by hand. These are teaching inputs, not evidence from a random experiment. The statistical experiments use actual draws from the generator.

\section{More than one answer from a function}
A \emph{named tuple} groups a few values and gives each a name. The model's transition function returns three probabilities. A result with fields \code{up}, \code{down}, and \code{stay} lets us write \code{probabilities.up}. The dot asks for the value stored under that field name. This is easier to read than remembering that the first entry means an increase.

\section{Files, libraries, and the project}
\code{include} reads and evaluates a Julia source file. A \code{module} gives a related set of functions its own name. The supplied source is called \code{MoranModel}; after including it, \code{using .MoranModel} makes its exported functions available. The dot means the module was defined locally rather than installed as an external package.

The project uses standard libraries supplied with Julia. Its \code{Project.toml} identifies these dependencies; \code{Manifest.toml} records the resolved environment. A TOML file is a configuration file, not a Julia program. The README gives the exact commands, so you do not need to learn configuration syntax to begin.

\begin{checkpoint}
\textbf{Try it.} What are the values of \code{i}, \code{N-i}, and \code{i/N} after \code{N = 10} and \code{i = 3}? Is \code{i == N} true?\\[3pt]
\textbf{Answer.} The count is three, the other type has seven individuals, and the frequency is \(0.3\). The equality test is false: three is not ten. Nothing in that test changes either variable.
\end{checkpoint}

\begin{checkpoint}
\textbf{Try it.} Why can the first recorded state in an array be numbered one even though it represents event zero?\\[3pt]
\textbf{Answer.} Array position and biological event number are different labels. Julia numbers positions from one. We choose to record the population before any event, so that state occupies position one and has event number zero.
\end{checkpoint}
